\conclusion % Структурный элемент: ЗАКЛЮЧЕНИЕ

В ходе выполнения выпускной квалификационной работы была решена
актуальная задача разработки методики верификации соответствия
TLA+ спецификации и исходного кода распределённой системы.

В исследовательской части были рассмотрены модель распределённой
системы, основные свойства каналов связи, временные предположения и
модели отказов, существенные для построения отказоустойчивых
распределённых приложений. Выполнен обзор существующих алгоритмов
консенсуса и показано, что для разрабатываемой системы распределённых
вычислений алгоритм Raft является наиболее подходящим решением
благодаря наличию выделенного лидера, встроенной поддержке
реплицируемого журнала и более удобной практической реализации.
Также были проанализированы существующие подходы к верификации
соответствия формальной спецификации и программной реализации, после
чего была сформулирована математическая постановка задачи в терминах
системы переходов, трассы исполнения и предиката согласованности
шага.

В конструкторской части были обоснованы основные проектные решения,
включая выбор алгоритма консенсуса, языка программирования и
инструментов реализации. Спроектирована архитектура распределённой
системы вычислений, включающая кластер Raft-узлов, клиентскую часть и
вычислительные узлы, а также разработана схема взаимодействия
компонентов системы верификации. Кроме того, было обосновано решение
математической постановки задачи и показано, каким образом задача
проверки совместимости трассы с моделью может быть сведена к
практически реализуемой процедуре валидации трасс средствами TLA+ и
TLC.

В технологической части была разработана TLA+ спецификация алгоритма
Raft и выполнена её модельная проверка. На основе данной
спецификации реализована распределённая система вычислений,
поддерживающая согласованное хранение состояния, репликацию журнала,
обработку клиентских запросов и выполнение вычислительных задач на
рабочих узлах. Также была разработана библиотека инструментирования
реализации, позволяющая фиксировать логические шаги исполнения,
формировать трассы в формате NDJSON и подготавливать их к
последующей проверке.

Практическая ценность работы подтверждена результатами интеграции
системы верификации в программную реализацию. Показано, что
предложенный подход позволяет сопоставлять наблюдаемое исполнение
распределённой системы с TLA+ спецификацией и выявлять расхождения
между моделью и исходным кодом. На демонстрационном примере было
показано, что валидация трасс действительно позволяет обнаруживать
нетривиальные логические ошибки реализации, проявляющиеся лишь при
определённых сценариях взаимодействия и повторной доставке сообщений.
Тем самым подтверждена применимость предложенной методики для
контроля корректности распределённых программных систем.

Таким образом, цель работы достигнута: разработана методика
верификации соответствия TLA+ спецификации и исходного кода,
предложена её математическая основа, реализованы программные
средства поддержки данного подхода и выполнена апробация на
распределённой системе вычислений, использующей алгоритм Raft.

Полученные результаты могут быть использованы при разработке и
сопровождении распределённых систем, для которых критически важны
согласованность состояния, предсказуемость поведения при сбоях и
возможность раннего выявления ошибок реализации. Разработанная
методика не ограничивается только рассмотренным примером и может
рассматриваться как основа для дальнейшего расширения на другие
классы распределённых алгоритмов и программных систем.

В качестве направлений дальнейшего развития работы представляются
целесообразными расширение библиотеки инструментирования для
поддержки новых языков программирования, в первую очередь Lua, а
также интеграция предложенной методики в реальные программные
платформы, использующие данный язык. В частности, перспективным
направлением является интеграция системы верификации в СУБД
Tarantool, где Lua применяется как один из основных языков
разработки прикладной логики. Такое развитие позволит проверить
практическую применимость методики в промышленной среде и
расширить область её использования за пределы разработанного в
работе прототипа.
